\section{Ermittlung der auslösenden Faktoren}\label{sec:af}

Dieser Abschnitt dokumentiert das Verfahren zur jährlichen Ermittlung
der auslösenden Faktoren der Leistungen und der Sterblichkeit je
Beobachtungseinheit. In diesem Abschnitt werden keine Anlagen geführt.

\subsection{Allgemeines}\label{sec:af-allgemein}

Rechtsgrundlage für Beitragsanpassungen sind \S\,203 Abs.\,2 VVG
i.\,V.\,m. \S\,8b der AVB sowie \S\,155 VAG. Das Unternehmen vergleicht
jährlich je Beobachtungseinheit die erforderlichen mit den kalkulierten
Versicherungsleistungen (\S\,155 Abs.\,3 VAG, \S\,15 KVAV) sowie die
erforderlichen mit den kalkulierten Sterbewahrscheinlichkeiten
(\S\,155 Abs.\,4 VAG, \S\,16 KVAV). Beobachtungseinheiten für die
auslösenden Faktoren sind \textbf{Männer}, \textbf{Frauen} sowie
\textbf{Kinder und Jugendliche zusammengefasst}. Ergibt die
Gegenüberstellung der Leistungen eine Abweichung von mehr als
\textbf{5\,\%} (vertraglich in \S\,8b AVB festgelegter Schwellenwert
unterhalb des gesetzlichen Werts von 10\,\%) oder die
Gegenüberstellung der Sterbewahrscheinlichkeiten eine Abweichung von
mehr als \textbf{5\,\%}, werden alle Beiträge der Beobachtungseinheit
überprüft und, soweit die Abweichung nicht nur vorübergehend ist, mit
Zustimmung des Treuhänders angepasst. In diesem Abschnitt werden keine
Anlagen geführt.

\subsection{Tatsächliche Grundkopfschäden}\label{sec:af-gks}

Für jedes der letzten drei Beobachtungsjahre wird der tatsächliche
Grundkopfschaden gemäß Abschnitt~A der Anlage~2 zur KVAV ermittelt:
\begin{equation}
G^{(j)} = \frac{S^{(j)}}{\sum_{x} L^{(j)}_x \cdot k_x},
\qquad j = 1,2,3,
\label{eq:af-gks}
\end{equation}
wobei $S^{(j)}$ die abgegrenzten Schadenzahlungen der
Beobachtungseinheit im Beobachtungsjahr $j$ (ohne die auf
Nettorisikozuschläge entfallenden Leistungen, einschließlich der
geschlechtsunabhängig verteilten Leistungen für Schwangerschaft und
Mutterschaft), $L^{(j)}_x$ den abgegrenzten mittleren Bestand des
Alters $x$ und $k_x$ die rechnungsmäßigen Profilwerte gemäß
Anlage~\ref{sec:anlage233} bezeichnen. Bei zwischenzeitlichen Leistungsänderungen
werden die tatsächlichen Grundkopfschäden auf das aktuelle
Leistungsversprechen umgerechnet (\S\,15 Abs.\,2 KVAV); die
Umrechnungsfaktoren werden im Anpassungsfall dokumentiert.

\textbf{Verfahrensdetails (Festlegungen):} Die Bestände $L^{(j)}_x$
werden als mittlere Personenjahre des Beobachtungsjahres ohne
Elimination von Beginnjahren ermittelt (der Tarif ist für das
Neugeschäft geschlossen, Wartezeit- und Selektionseffekte sind
abgeklungen). Die Schäden $S^{(j)}$ werden nach dem Anfalljahrprinzip
abgegrenzt; der Schaden des jüngsten Beobachtungsjahres wird mit
Abwicklungsfaktoren aus der unternehmenseigenen
Abwicklungsstatistik auf den Endstand hochgerechnet. Für die
zusammengefasste Beobachtungseinheit der Kinder und Jugendlichen liegt
kein Altersprofil vor (natürliche Prämie); es wird $k_x := 1$ für alle
$x \le 20$ gesetzt (altersunabhängige Kopfschäden gemäß
Abschnitt~\ref{sec:kopfschaden-kinder}), sodass sich der tatsächliche
Grundkopfschaden gemäß Gleichung~\eqref{eq:af-gks} zum
bestandsgewichteten mittleren Kopfschaden
$G^{(j)} = S^{(j)} / \sum_x L^{(j)}_x$ vereinfacht. Als rechnungsmäßiger
Vergleichswert dient der bestandsgewichtete Mittelwert der
rechnungsmäßigen Kopfschäden
$G^{\mathrm{rech}} = 0{,}75\cdot\bar{K}_{\mathrm{Ki}}
+ 0{,}25\cdot\bar{K}_{\mathrm{Ju}}$, wobei $\bar{K}_{\mathrm{Ju}}$ das
bestandsgewichtete Mittel der männlichen und weiblichen Jugendlichen
bezeichnet.

\subsection{Auslösender Faktor der Leistungen}\label{sec:af-leistung}

Der extrapolierte Grundkopfschaden gemäß Abschnitt~B der Anlage~2 zur
KVAV (lineare Trendgerade durch die drei Beobachtungswerte,
ausgewertet 18~Monate nach dem Ende des letzten Beobachtungsjahres)
beträgt
\begin{equation}
G^{\mathrm{ext}} = \frac{3}{2}\cdot\bigl(G^{(3)} - G^{(1)}\bigr)
+ \frac{1}{3}\cdot\bigl(G^{(1)} + G^{(2)} + G^{(3)}\bigr),
\label{eq:af-ext}
\end{equation}
wobei $G^{(1)}, G^{(2)}, G^{(3)}$ die (umgerechneten) tatsächlichen
Grundkopfschäden in zeitlicher Reihenfolge bezeichnen. Mit dem
mittleren Bestand $L_x$ des aktuellen Geschäftsjahres ergeben sich die
erforderlichen und die kalkulierten Versicherungsleistungen
\begin{equation}
S^{\mathrm{erf}} = G^{\mathrm{ext}}\cdot\sum_x L_x\,k_x,
\qquad
S^{\mathrm{kalk}} = G^{\mathrm{rech}}\cdot\sum_x L_x\,k_x,
\label
{eq:af-serf}
\end{equation}
wobei $G^{\mathrm{rech}}$ den rechnungsmäßigen Grundkopfschaden gemäß
Gleichung~\eqref{eq:gks} bezeichnet. Der auslösende Faktor der
Leistungen ist
\begin{equation}
AF^{L} = \frac{S^{\mathrm{erf}}}{S^{\mathrm{kalk}}}
       = \frac{G^{\mathrm{ext}}}{G^{\mathrm{rech}}};
\label{eq:af-l}
\end{equation}
eine Überprüfungspflicht besteht, falls
$\bigl|1 - AF^{L}\bigr| > 0{,}05$.

\subsection{Auslösender Faktor der Sterblichkeit}\label{sec:af-sterb}

Gemäß \S\,16 KVAV werden die Leistungsbarwerte \textbf{ohne
Berücksichtigung der Stornowahrscheinlichkeiten} mit dem
rechnungsmäßigen Zins und den rechnungsmäßigen Kopfschäden verglichen:
\begin{equation}
A^{\mathrm{kalk}}_x = \sum_{t=0}^{110-x} v^t\, K_{x+t}
\prod_{s=0}^{t-1}\bigl(1-q_{x+s}\bigr),
\qquad
A^{\mathrm{erf}}_x = \sum_{t=0}^{110-x} v^t\, K_{x+t}
\prod_{s=0}^{t-1}\bigl(1-q^{\mathrm{BaFin}}_{x+s}\bigr),
\label{eq:af-barwerte}
\end{equation}
wobei $q_{x}$ die rechnungsmäßigen Sterbewahrscheinlichkeiten gemäß
Abschnitt~\ref{sec:sterblichkeit} und $q^{\mathrm{BaFin}}_{x}$ die
Werte der zuletzt von der BaFin veröffentlichten PKV-Sterbetafel
bezeichnen. Der auslösende Faktor der Sterblichkeit ist das Maximum
der über drei Altersbänder gemittelten Barwertquotienten:
\begin{equation}
AF^{S} = \max\left\{
\frac{1}{25}\sum_{x=21}^{45}\frac{A^{\mathrm{erf}}_x}{A^{\mathrm{kalk}}_x};\;
\frac{1}{25}\sum_{x=46}^{70}\frac{A^{\mathrm{erf}}_x}{A^{\mathrm{kalk}}_x};\;
\frac{1}{25}\sum_{x=71}^{95}\frac{A^{\mathrm{erf}}_x}{A^{\mathrm{kalk}}_x}
\right\};
\label{eq:af-s}
\end{equation}
eine Überprüfungspflicht besteht, falls
$\bigl|1 - AF^{S}\bigr| > 0{,}05$ (\S\,155 Abs.\,4 VAG).

\subsection{Fristen und Verfahren}\label{sec:af-verfahren}

Die auslösenden Faktoren werden jährlich bis spätestens vier Monate
nach Ende des Beobachtungsjahres berechnet und dem Treuhänder sowie
der Aufsichtsbehörde vorgelegt (\S\,17 KVAV). Springt mindestens ein
auslösender Faktor einer Beobachtungseinheit an, werden alle
Rechnungsgrundlagen dieser Beobachtungseinheit überprüft und die
Beiträge gemäß Abschnitt~\ref{sec:bap} neu berechnet; der
Sicherheitszuschlag darf dabei nicht erhöht werden, soweit die
Anpassung hierauf beruht. Soweit Leistungen bei der Erst- oder einer
Neukalkulation erkennbar unzureichend kalkuliert waren, erfolgt
insoweit keine Anpassung (\S\,155 Abs.\,3 VAG).
